% GENERATED FILE. Edit canonical lessons, then run npm run generate:books.
% canonical-source-hash: fnv1a64:98f40e5130eedafe
% canonical-lessons: PA-C243-pauna, PA-C243-gholna, PA-C243-ralauna, PA-C243-chhanna, PA-C243-chakkhna

\chapter{Mixing and Tasting}
\label{ch:pa-a2-mixing-and-tasting}
\hlchaptermodality{243}

\begin{chapteropening}
\textbf{By the end of this chapter:} \emph{I can say how I mix and taste food.}

Complete the last lesson of chapter 243: I can say how I mix and taste food.
\end{chapteropening}

\section[\texorpdfstring{\pa{ਪਾਉਣਾ} (pāuṇā)}{ਪਾਉਣਾ (pāuṇā)}]{\pa{ਪਾਉਣਾ} (pāuṇā) — to pour in; to put in}

\label{lesson:PA-C243-pauna}



\begin{quote}
Before the new one: say the Punjabi for to sprinkle, then the Punjabi for to get burnt.
\end{quote}

\subsection*{You'll want to know: \pa{ਪਾਉਣਾ}}
\textbf{\pa{ਪਾਉਣਾ}} — \emph{pāuṇā} — \textquotedblleft{}to pour in; to put in\textquotedblright{}.

\begin{quote}
\pa{ਮੈਂ} \pa{ਚਾਹ} \pa{ਵਿੱਚ} \pa{ਦੁੱਧ} \pa{ਪਾਉਂਦੀ} \pa{ਹਾਂ।}
\end{quote}

\begin{quote}
\emph{maiṁ chāh vich dudh pāundī hāṁ.}
\end{quote}

\begin{quote}
I pour milk into the tea.
\end{quote}

Past (I …): \pa{ਮੈਂ} \pa{ਦੁੱਧ} \pa{ਪਾਇਆ}, \emph{maiṁ dudh pāiā}.

\begin{cousinweb}
\textbf{\pa{ਪਾਉਣਾ}} (\emph{pāuṇā}) is inherited from Sanskrit \emph{pātayati}, \textquotedblleft{}throws, pours out\textquotedblright{}; CDIAL lists Punjabi \emph{pāuṇā}, \textquotedblleft{}to pour\textquotedblright{} (CDIAL 8053).
\end{cousinweb}

\begin{grammarlens}[title={What you've built}]
One more word for this chapter.
\end{grammarlens}

\subsection*{Your turn}
\begin{itemize}
\raggedright
  \item \emph{Say it:} \emph{pāuṇā}
  \item \emph{Say it:} \emph{pāuṇā}, once more
  \item \emph{Say it:} \emph{maiṁ chāh vich dudh pāundī hāṁ.} — I pour milk into the tea.
  \item \emph{Recall:} say the Punjabi for to sprinkle, then the Punjabi for to get burnt, then say \emph{pāuṇā} again
\end{itemize}

\begin{tcolorbox}[breakable,colback=teal!4,colframe=teal!35!black,title={Before you move on}]
What does \pa{ਪਾਉਣਾ} mean? (\textquotedblleft{}To pour in; to put in\textquotedblright{}.) Say it once more.
\end{tcolorbox}
\section[\texorpdfstring{\pa{ਘੋਲਣਾ} (gholṇā)}{ਘੋਲਣਾ (gholṇā)}]{\pa{ਘੋਲਣਾ} (gholṇā) — to dissolve; to stir in}

\label{lesson:PA-C243-gholna}



\begin{quote}
Before the new one: say the Punjabi for to get burnt, then the Punjabi for to pour in.
\end{quote}

\subsection*{You'll want to know: \pa{ਘੋਲਣਾ}}
\textbf{\pa{ਘੋਲਣਾ}} — \emph{gholṇā} — \textquotedblleft{}to dissolve; to stir in\textquotedblright{}.

\begin{quote}
\pa{ਮੈਂ} \pa{ਪਾਣੀ} \pa{ਵਿੱਚ} \pa{ਖੰਡ} \pa{ਘੋਲਦਾ} \pa{ਹਾਂ।}
\end{quote}

\begin{quote}
\emph{maiṁ pāṇī vich khaṇḍ gholdā hāṁ.}
\end{quote}

\begin{quote}
I dissolve sugar in water.
\end{quote}

Past (I …): \pa{ਮੈਂ} \pa{ਖੰਡ} \pa{ਘੋਲੀ}, \emph{maiṁ khaṇḍ gholī}.

\begin{cousinweb}
\textbf{\pa{ਘੋਲਣਾ}} (\emph{gholṇā}) is inherited from Sanskrit \emph{ghōlayati}, \textquotedblleft{}mixes, stirs together\textquotedblright{}; CDIAL lists Punjabi \emph{gholaṇā}, \textquotedblleft{}to stir, mix, dissolve\textquotedblright{} (CDIAL 4526).
\end{cousinweb}

\begin{grammarlens}[title={What you've built}]
One more word for this chapter.
\end{grammarlens}

\subsection*{Your turn}
\begin{itemize}
\raggedright
  \item \emph{Say it:} \emph{gholṇā}
  \item \emph{Say it:} \emph{gholṇā}, once more
  \item \emph{Say it:} \emph{maiṁ pāṇī vich khaṇḍ gholdā hāṁ.} — I dissolve sugar in water.
  \item \emph{Recall:} say the Punjabi for to get burnt, then the Punjabi for to pour in, then say \emph{gholṇā} again
  \item \emph{Write it:} copy \textbf{\pa{ਘੋਲਣਾ}} once, checking the whole word
\end{itemize}

\begin{tcolorbox}[breakable,colback=teal!4,colframe=teal!35!black,title={Before you move on}]
What does \pa{ਘੋਲਣਾ} mean? (\textquotedblleft{}To dissolve; to stir in\textquotedblright{}.) Say it once more.
\end{tcolorbox}
\section[\texorpdfstring{\pa{ਰਲਾਉਣਾ} (ralāuṇā)}{ਰਲਾਉਣਾ (ralāuṇā)}]{\pa{ਰਲਾਉਣਾ} (ralāuṇā) — to mix}

\label{lesson:PA-C243-ralauna}



\begin{quote}
Before the new one: say the Punjabi for to pour in, then the Punjabi for to dissolve.
\end{quote}

\subsection*{You'll want to know: \pa{ਰਲਾਉਣਾ}}
\textbf{\pa{ਰਲਾਉਣਾ}} — \emph{ralāuṇā} — \textquotedblleft{}to mix\textquotedblright{}.

\begin{quote}
\pa{ਮੈਂ} \pa{ਦਾਲ} \pa{ਵਿੱਚ} \pa{ਮਸਾਲਾ} \pa{ਰਲਾਉਂਦੀ} \pa{ਹਾਂ।}
\end{quote}

\begin{quote}
\emph{maiṁ dāl vich masālā ralāundī hāṁ.}
\end{quote}

\begin{quote}
I mix spices into the dal.
\end{quote}

Past (I …): \pa{ਮੈਂ} \pa{ਮਸਾਲਾ} \pa{ਰਲਾਇਆ}, \emph{maiṁ masālā ralāiā}.

\begin{cousinweb}
\textbf{\pa{ਰਲਾਉਣਾ}} (\emph{ralāuṇā}) is inherited from a reconstructed \emph{*ral}, \textquotedblleft{}meet with, join\textquotedblright{}; CDIAL lists Punjabi \emph{ralāuṇā}, \textquotedblleft{}to mix, join\textquotedblright{} (CDIAL 10640).
\end{cousinweb}

\begin{grammarlens}[title={What you've built}]
One more word for this chapter.
\end{grammarlens}

\subsection*{Your turn}
\begin{itemize}
\raggedright
  \item \emph{Say it:} \emph{ralāuṇā}
  \item \emph{Say it:} \emph{ralāuṇā}, once more
  \item \emph{Say it:} \emph{maiṁ dāl vich masālā ralāundī hāṁ.} — I mix spices into the dal.
  \item \emph{Recall:} say the Punjabi for to pour in, then the Punjabi for to dissolve, then say \emph{ralāuṇā} again
  \item \emph{Write it:} copy \textbf{\pa{ਰਲਾਉਣਾ}} once, checking the whole word
\end{itemize}

\begin{tcolorbox}[breakable,colback=teal!4,colframe=teal!35!black,title={Before you move on}]
What does \pa{ਰਲਾਉਣਾ} mean? (\textquotedblleft{}To mix\textquotedblright{}.) Say it once more.
\end{tcolorbox}
\section[\texorpdfstring{\pa{ਛਾਣਨਾ} (chhāṇnā)}{ਛਾਣਨਾ (chhāṇnā)}]{\pa{ਛਾਣਨਾ} (chhāṇnā) — to sift}

\label{lesson:PA-C243-chhanna}



\begin{quote}
Before the new one: say the Punjabi for to dissolve, then the Punjabi for to mix.
\end{quote}

\subsection*{You'll want to know: \pa{ਛਾਣਨਾ}}
\textbf{\pa{ਛਾਣਨਾ}} — \emph{chhāṇnā} — \textquotedblleft{}to sift\textquotedblright{}.

\begin{quote}
\pa{ਮੈਂ} \pa{ਆਟਾ} \pa{ਛਾਣਦੀ} \pa{ਹਾਂ।}
\end{quote}

\begin{quote}
\emph{maiṁ āṭā chhāṇdī hāṁ.}
\end{quote}

\begin{quote}
I sift the flour.
\end{quote}

Past (I …): \pa{ਮੈਂ} \pa{ਆਟਾ} \pa{ਛਾਣਿਆ}, \emph{maiṁ āṭā chhāṇiā}.

\begin{cousinweb}
\textbf{\pa{ਛਾਣਨਾ}} (\emph{chhāṇnā}) is inherited from a reconstructed \emph{*kṣaṇati}, \textquotedblleft{}is sifted, is strained\textquotedblright{}; CDIAL lists Punjabi \emph{chāṇnā}, \textquotedblleft{}to sift, strain, choose\textquotedblright{} (CDIAL 3643).
\end{cousinweb}

\begin{grammarlens}[title={What you've built}]
One more word for this chapter.
\end{grammarlens}

\subsection*{Your turn}
\begin{itemize}
\raggedright
  \item \emph{Say it:} \emph{chhāṇnā}
  \item \emph{Say it:} \emph{chhāṇnā}, once more
  \item \emph{Say it:} \emph{maiṁ āṭā chhāṇdī hāṁ.} — I sift the flour.
  \item \emph{Recall:} say the Punjabi for to dissolve, then the Punjabi for to mix, then say \emph{chhāṇnā} again
  \item \emph{Write it:} copy \textbf{\pa{ਛਾਣਨਾ}} once, checking the whole word
\end{itemize}

\begin{tcolorbox}[breakable,colback=teal!4,colframe=teal!35!black,title={Before you move on}]
What does \pa{ਛਾਣਨਾ} mean? (\textquotedblleft{}To sift\textquotedblright{}.) Say it once more.
\end{tcolorbox}
\section[\texorpdfstring{\pa{ਚੱਖਣਾ} (chakkhṇā)}{ਚੱਖਣਾ (chakkhṇā)}]{\pa{ਚੱਖਣਾ} (chakkhṇā) — to taste}

\label{lesson:PA-C243-chakkhna}



\begin{quote}
Before the new one: say the Punjabi for to mix, then the Punjabi for to sift.
\end{quote}

\subsection*{You'll want to know: \pa{ਚੱਖਣਾ}}
\textbf{\pa{ਚੱਖਣਾ}} — \emph{chakkhṇā} — \textquotedblleft{}to taste\textquotedblright{}.

\begin{quote}
\pa{ਕੀ} \pa{ਤੁਸੀਂ} \pa{ਖੀਰ} \pa{ਚੱਖੋਗੇ}?
\end{quote}

\begin{quote}
\emph{kī tusī̃ khīr chakkhoge?}
\end{quote}

\begin{quote}
Will you taste the rice pudding?
\end{quote}

Past (I …): \pa{ਮੈਂ} \pa{ਖੀਰ} \pa{ਚੱਖੀ}, \emph{maiṁ khīr chakkhī}.

\begin{cousinweb}
\textbf{\pa{ਚੱਖਣਾ}} (\emph{chakkhṇā}) is inherited from a reconstructed \emph{*cakṣati}, \textquotedblleft{}tastes\textquotedblright{}; CDIAL lists Punjabi \emph{cakkhṇā}, \textquotedblleft{}to taste\textquotedblright{} (CDIAL 4557).
\end{cousinweb}

\begin{grammarlens}[title={What you've built}]
One more word for this chapter.
\end{grammarlens}

\subsection*{Your turn}
\begin{itemize}
\raggedright
  \item \emph{Say it:} \emph{chakkhṇā}
  \item \emph{Say it:} \emph{chakkhṇā}, once more
  \item \emph{Say it:} \emph{kī tusī̃ khīr chakkhoge?} — Will you taste the rice pudding?
  \item \emph{Recall:} say the Punjabi for to mix, then the Punjabi for to sift, then say \emph{chakkhṇā} again
  \item \emph{Write it:} copy \textbf{\pa{ਚੱਖਣਾ}} once, checking the whole word
\end{itemize}

\begin{tcolorbox}[breakable,colback=teal!4,colframe=teal!35!black,title={Before you move on}]
What does \pa{ਚੱਖਣਾ} mean? (\textquotedblleft{}To taste\textquotedblright{}.) Say it once more.
\end{tcolorbox}
